% Part: computability
% Chapter: computability-theory
% Section: ce-closed-cup-cap

\documentclass[../../../include/open-logic-section]{subfiles}

\begin{document}

\olfileid[pt]{cmp}{thy}{clo}

\olsection[União e interseção de conjuntos c.e.]{Conjuntos
  computavelmente enumeráveis são fechados sob união e interseção}

O teorema a seguir apresenta algumas propriedades de fechamento do
conjunto dos conjuntos computavelmente enumeráveis.

\begin{thm}
Suponha que $A$ e $B$ sejam computavelmente enumeráveis. Então
$A \cap B$ e $A \cup B$ também o são.
\end{thm}

\begin{proof}
A \olref[eqc]{thm:ce-equiv} nos permite usar várias caracterizações
dos conjuntos computavelmente enumeráveis. Para ilustrar, apresentaremos
algumas demonstrações diferentes.

Para a primeira demonstração, suponha que $A$ seja enumerado por uma
função computável~$f$ e que $B$ seja enumerado por uma função
computável~$g$. Seja
\begin{align*}
h(x) & = \umin{y}{(f(y) = x \lor g(y) = x)} \text{ e}\\
j(x) & = \umin{y}{(f((y)_0) = x \land g((y)_1) = x)}.
\end{align*}
Então $A \cup B$ é o domínio de $h$, e $A \cap B$ é o domínio de~$j$.

\begin{explain}
Em termos computacionais, ocorre o seguinte: dados procedimentos que
enumeram $A$ e $B$, podemos semidecidir se um elemento $x$ pertence a
$A \cup B$ procurando $x$ em qualquer uma das enumerações; e podemos
semidecidir se um elemento $x$ pertence a $A \cap B$ procurando $x$
nas duas enumerações ao mesmo tempo.
\end{explain}

Para a segunda demonstração, suponha novamente que $A$ seja enumerado
por~$f$ e $B$ por~$g$. Seja
\[
k(x) = \begin{cases}
f(x/2) & \text{se $x$ é par} \\
g((x-1)/2) & \text{se $x$ é ímpar.}
\end{cases}
\]
Então $k$ enumera $A \cup B$; a ideia é que $k$ simplesmente alterna
entre as enumerações fornecidas por $f$ e~$g$. Enumerar $A \cap B$ é
mais trabalhoso. Se $A \cap B$ for vazio, será trivialmente
computavelmente enumerável. Caso contrário, seja $c$ um elemento
qualquer de $A \cap B$ e defina $l$ por
\[
l(x) = \begin{cases}
f((x)_0) & \text{se $f((x)_0) = g((x)_1)$} \\
c & \text{caso contrário.}
\end{cases}
\]
Em termos computacionais, $l$ percorre pares de elementos nas
enumerações de $f$ e $g$ e produz como saída cada coincidência que
encontra; caso contrário, apenas fica inativa, produzindo $c$ como
saída.

Para a última demonstração, suponha que $A$ seja o \emph{domínio} da
função parcial $m(x)$ e que $B$ seja o domínio da função parcial
$n(x)$. Então $A \cap B$ é o domínio da função parcial $m(x) + n(x)$.

\begin{explain}
Em termos computacionais, se $A$ é o conjunto dos valores para os quais
$m$ para e $B$ é o conjunto dos valores para os quais $n$ para,
$A \cap B$ é o conjunto dos valores para os quais ambos os procedimentos
param.
\end{explain}

Expressar $A \cup B$ como um conjunto de valores de parada é mais
difícil, pois é preciso simular $m$ e $n$ em paralelo. Seja $d$ um
índice de $m$ e seja $e$ um índice de $n$; em outras palavras,
$m = \cfind{d}$ e $n = \cfind{e}$. Então $A \cup B$ é o domínio da
função
\[
p(x) = \umin{y}{(T(d,x,y) \lor T(e,x,y))}.
\]
\begin{explain}
Em termos computacionais, na entrada $x$, $p$ procura uma computação
que pare de $m$ ou uma computação que pare de $n$, e para se encontrar
qualquer uma delas.
\end{explain}
\end{proof}

\end{document}

